\documentclass[11pt]{amsart}
\usepackage{geometry}                % See geometry.pdf to learn the layout options. There are lots.
\geometry{letterpaper}                   % ... or a4paper or a5paper or ... 
%\geometry{landscape}                % Activate for for rotated page geometry
%\usepackage[parfill]{parskip}    % Activate to begin paragraphs with an empty line rather than an indent
\usepackage{graphicx}
\usepackage{amssymb}
\usepackage{epstopdf}
\usepackage{color}
\usepackage{colortbl}
\DeclareGraphicsRule{.tif}{png}{.png}{`convert #1 `dirname #1`/`basename #1 .tif`.png}

\title{PS 3.25}
%\date{}                                           % Activate to display a given date or no date

\begin{document}
\maketitle
%\section{}
%\subsection{}

\noindent En una escuela en la que se tienen registros con las características físicas de los alumnos, se desea conocer la lista de los alumnos con aptitudes para practicar básquet. Haga un diagrama de flujo que obtenga lo siguiente:\\

\noindent a) Lista de alumnas con aptitudes físicas para jugar al básquet.\\

\noindent \quad $Requerimientos : ALTURA > 1. 73 \quad y \quad 50 < PESO < 90$.\\

\noindent b) Porcentaje de alumnas con estas aptitudes de la población estudiantil femenina.\\

\noindent c) Lista de alumnos con aptitudes físicas para jugar al básquet.\\

\noindent \quad $Requerimientos : ALTURA > 1.83  \quad y  \quad 73 < PESO < 110$.\\

\noindent d) Porcentaje de alumnos con estas aptitudes de la población estudiantil masculina.\\

\noindent Por cada alumno se ingresa su NOMBRE, SEXO, EDAD, PESO y ALTURA.\\

\noindent Datos: $NOM_{1}, SEX_{1}, EDAD_{1}, PESO_{1}, ALT_{1}, NOM_{2}, SEX_{2}, EDAD_{2}, PESO_{2}, ALT_{2}> . . ., “ X ”, “ X ”,- 1.-1,-$\\

\noindent Donde:\\

\noindent \quad $NOM_i$ es una variable de tipo cadena de caracteres que representa el NOMBRE del alumno i.\\

\noindent \quad $SEX_i$ es una variable de tipo carácter que expresa el SEXO del alumno i. Se ingresa “ F ” para mujer y “ M ” para hombre.\\

\noindent \quad $EDAD_i$ es una variable de tipo entero que representa la EDAD del alumno i.\\

\noindent \quad $PESO_i$ ees una variable de tipo entero que representa el PESO del alumno i.\\

\noindent \quad $ALT_i$ es una variable de tipo real que representa la ALTURA del alumno i.\\

\end{document}   